%%%%%%%%%%%%%%%%%%%%%%%%%%%%%%%%%%%%%%%%%%%%%%%%%%%%%%%%%%%%%%%%%%%%%%%%%%%%%%%%
%% 02_Projeto_Logico_Rede.tex
%% Projeto Lógico da Rede, Endereçamento IPv4, VLANs e Políticas de Roteamento
%%%%%%%%%%%%%%%%%%%%%%%%%%%%%%%%%%%%%%%%%%%%%%%%%%%%%%%%%%%%%%%%%%%%%%%%%%%%%%%%

\chapter{Projeto Lógico da Rede de Computadores}
\label{cap:projeto_logico}

\section{Arquitetura e Modelo Topológico Hierárquico}

O projeto lógico adota a arquitetura de rede hierárquica dividida em camadas de borda, segurança perimetral, distribuição e acesso\footnote{Kurose, J. F.; Ross, K. W. Redes de Computadores e a Internet: Uma Abordagem Top-Down. 8. ed. Pearson, 2021.}. A divisão funcional isola domínios de \textit{broadcast}, simplifica a resolução de falhas operacionais e padroniza a aplicação de políticas de segurança.

A topologia lógica emulada no simulador Cisco Packet Tracer estrutura-se em três blocos operacionais delimitados, conforme apresentado na Figura~\ref{fig:topologia_logica}:
\begin{enumerate}
    \item \textbf{Bloco Externo (Infraestrutura do Provedor)}: Composto pelo roteador \texttt{ISP-UFS} (Cisco 2911), representando a operadora acadêmica da Universidade Federal de Sergipe. O dispositivo fornece o enlace de trânsito WAN público (\texttt{200.17.0.0/30}) e disponibiliza o endereço público \texttt{8.8.8.8/32} em sua interface lógica \textit{Loopback0}, viabilizando testes de resolução de nomes (DNS) e verificação de saída para a rede externa;
    \item \textbf{Bloco CT12 (Laboratório de Geoprocessamento e Servidores, Sala 201)}: Concentra os ativos de borda, segurança física, roteamento de núcleo e servidores de aplicação e dados:
    \begin{itemize}
        \item \textbf{Roteador de Borda (\texttt{router\_border})}: Modelo corporativo Cisco ISR 4321, dotado de interface WAN interligada ao \texttt{ISP-UFS} e interface interna direcionada ao \textit{firewall}. Executa a tradução de endereços por sobrecarga (\textit{NAT Overload}), garantindo saída transparente para toda a rede interna sob um único IP público;
        \item \textbf{Firewall Físico (\texttt{Firewall})}: Modelo Cisco ASA 5506-X, posicionado estrategicamente entre o roteador de borda e o roteador de distribuição interna. Cria duas zonas de segurança com níveis de confiança distintos (\textit{outside} com nível 0 e \textit{inside} com nível 100), inspecionando fluxos de estado via \textit{Modular Policy Framework} (MPF) para os protocolos ICMP, HTTP e DNS;
        \item \textbf{Roteador Interno (\texttt{router\_interno})}: Modelo Cisco 2911, responsável pelo roteamento entre as sub-redes virtuais por meio de enlaces troncais 802.1Q (\textit{Router-on-a-Stick}), aplicação de Listas de Controle de Acesso (ACL) estendidas e retransmissão de solicitações DHCP (\textit{DHCP Relay Agent} via \texttt{ip helper-address}) direcionadas ao servidor de infraestrutura;
        \item \textbf{Switch Central CT-12 (\texttt{SW-CT12})}: Modelo Cisco Catalyst 2960-24TT, comutador gerenciável que agrega os servidores dedicados da VLAN 20 (\texttt{SRV-NEXTGIS}, \texttt{SRV-POSTGIS} e \texttt{SRV-DHCP}), o ponto de acesso sem fio \texttt{AP-CT12}, a estação administrativa \texttt{WS-ADMIN}, as estações de pesquisa locais \texttt{WS-CT12-01} e \texttt{WS-CT12-02}, e provê o enlace tronco inter-edifícios;
    \end{itemize}
    \item \textbf{Bloco CT09 (Sala dos Pesquisadores, Sala 101)}:
    \begin{itemize}
        \item \textbf{Switch de Acesso CT-09 (\texttt{SW-CT09})}: Modelo Cisco Catalyst 2960-24TT, interligado ao \texttt{SW-CT12} por enlace tronco Gigabit Ethernet (fibra óptica monomodo no projeto físico). Conecta as estações de trabalho \texttt{WS-CT09-01} e \texttt{WS-CT09-02} (VLAN 30) e o ponto de acesso sem fio \texttt{AP-CT09} (VLAN 40).
    \end{itemize}
\end{enumerate}

\incluirdiagrama{Imagens/topologia_logica.png}{Topologia lógica integrada no Cisco Packet Tracer com Provedor ISP-UFS, Firewall Cisco ASA 5506-X e segmentação de VLANs entre CT-12 e CT-09.}{fig:topologia_logica}

\section{Plano de Endereçamento IPv4 e Segmentação em VLANs}

A rede local emprega blocos de endereçamento IPv4 privado (RFC 1918) e segmentos públicos para o trânsito com o provedor institucional. O particionamento compreende quatro Redes Locais Virtuais (VLANs) de serviço e três sub-redes de trânsito ponto a ponto (/30), assegurando contenção de falhas e proteção perimetral.

A Tabela~\ref{tab:enderecamento_vlan} apresenta o mapeamento de sub-redes, máscaras CIDR e portas de conexão padrão.

\begin{table}[htbp]
\centering
\small
\caption{Plano de endereçamento IP, segmentação de VLANs e enlaces de trânsito.}
\label{tab:enderecamento_vlan}
\begin{tabularx}{\textwidth}{|c|l|Y|l|r|}
\hline
\rowcolor{tableheaderbeige}
\textbf{VLAN / Enlace} & \textbf{Nome da Rede} & \textbf{Faixa de Rede (CIDR)} & \textbf{Gateway Padrão} & \textbf{Capacidade} \\ \hline
10 & Gerência de Ativos & \texttt{192.168.10.0/24} & \texttt{192.168.10.1} & 254 hosts \\ \hline
20 & Servidores e DMZ & \texttt{192.168.20.0/24} & \texttt{192.168.20.1} & 254 hosts \\ \hline
30 & Estações de Pesquisa & \texttt{192.168.30.0/24} & \texttt{192.168.30.1} & 254 hosts \\ \hline
40 & Rede Sem Fio (Wi-Fi) & \texttt{192.168.40.0/24} & \texttt{192.168.40.1} & 254 hosts \\ \hline
-- & Enlace WAN ISP-Borda & \texttt{200.17.0.0/30} & \texttt{200.17.0.1} & 2 hosts \\ \hline
-- & Enlace WAN-Firewall & \texttt{10.0.0.0/30} & \texttt{10.0.0.1} & 2 hosts \\ \hline
-- & Enlace Firewall-Núcleo & \texttt{10.0.1.0/30} & \texttt{10.0.1.1} & 2 hosts \\ \hline
-- & DNS/Internet Simulado & \texttt{8.8.8.8/32} & \texttt{8.8.8.8} & 1 host \\ \hline
\end{tabularx}
\end{table}

A Tabela~\ref{tab:enderecamento_dispositivos} sintetiza a atribuição de endereços IP nominais aos dispositivos ativos e nós terminais ilustrados na topologia lógica.

\begin{table}[htbp]
\centering
\small
\caption{Mapeamento de interfaces e atribuição de endereços aos dispositivos da topologia.}
\label{tab:enderecamento_dispositivos}
\begin{tabularx}{\textwidth}{|l|l|l|Y|}
\hline
\rowcolor{tableheaderbeige}
\textbf{Dispositivo} & \textbf{Interface} & \textbf{Endereço IP / Máscara} & \textbf{Função Operacional} \\ \hline
\texttt{ISP-UFS} & \texttt{GigabitEthernet0/0} & \texttt{200.17.0.1/30} & Gateway WAN da operadora acadêmica \\ \hline
\texttt{ISP-UFS} & \texttt{Loopback0} & \texttt{8.8.8.8/32} & Servidor DNS e destino público simulado \\ \hline
\texttt{router\_border} & \texttt{GigabitEthernet0/0/1} & \texttt{200.17.0.2/30} & Interface WAN com NAT Overload \\ \hline
\texttt{router\_border} & \texttt{GigabitEthernet0/0/0} & \texttt{10.0.0.1/30} & Ponto a ponto para o Firewall \\ \hline
\texttt{Firewall} & \texttt{GigabitEthernet1/1} & \texttt{10.0.0.2/30} & Interface \textit{outside} (Nível de segurança 0) \\ \hline
\texttt{Firewall} & \texttt{GigabitEthernet1/2} & \texttt{10.0.1.1/30} & Interface \textit{inside} (Nível de segurança 100) \\ \hline
\texttt{router\_interno} & \texttt{GigabitEthernet0/1} & \texttt{10.0.1.2/30} & Ponto a ponto para o Firewall (Rota default) \\ \hline
\texttt{router\_interno} & \texttt{GigabitEthernet0/0.10} & \texttt{192.168.10.1/24} & Gateway da VLAN 10 (Gerência) \\ \hline
\texttt{router\_interno} & \texttt{GigabitEthernet0/0.20} & \texttt{192.168.20.1/24} & Gateway da VLAN 20 (Servidores) \\ \hline
\texttt{router\_interno} & \texttt{GigabitEthernet0/0.30} & \texttt{192.168.30.1/24} & Gateway da VLAN 30 (\textit{DHCP Relay}) \\ \hline
\texttt{router\_interno} & \texttt{GigabitEthernet0/0.40} & \texttt{192.168.40.1/24} & Gateway da VLAN 40 (\textit{DHCP Relay}) \\ \hline
\texttt{SW-CT12} & \texttt{VLAN 10 (SVI)} & \texttt{192.168.10.2/24} & Gerenciamento do switch central CT-12 \\ \hline
\texttt{SW-CT09} & \texttt{VLAN 10 (SVI)} & \texttt{192.168.10.3/24} & Gerenciamento do switch de acesso CT-09 \\ \hline
\texttt{SRV-NEXTGIS} & \texttt{FastEthernet0} & \texttt{192.168.20.10/24} & Servidor da aplicação NextGIS Web (Portas 80/443) \\ \hline
\texttt{SRV-POSTGIS} & \texttt{FastEthernet0} & \texttt{192.168.20.20/24} & Servidor do banco PostGIS (Porta 5432) \\ \hline
\texttt{SRV-DHCP} & \texttt{FastEthernet0} & \texttt{192.168.20.30/24} & Servidor DHCP dedicado para VLANs 30 e 40 \\ \hline
\texttt{WS-ADMIN} & \texttt{FastEthernet0} & \texttt{192.168.30.10/24} & Estação autorizada de administração técnica \\ \hline
\texttt{WS-CT12-01/02} & \texttt{FastEthernet0} & Dinâmico (DHCP) & Estações de pesquisa do Prédio CT-12 \\ \hline
\texttt{WS-CT09-01/02} & \texttt{FastEthernet0} & Dinâmico (DHCP) & Estações de pesquisa do Prédio CT-09 \\ \hline
\texttt{AP-CT12/09} & \texttt{Ethernet0} & \texttt{192.168.10.11/12} & Pontos de acesso Wi-Fi 6 corporativos \\ \hline
\end{tabularx}
\end{table}

\section{Roteamento Inter-VLAN e Enlaces de Trânsito Ponto a Ponto}

O roteamento entre as sub-redes virtuais ocorre no roteador Cisco 2911 (\texttt{router\_interno}) através de enlaces troncais 802.1Q sobre a interface física Gigabit Ethernet \texttt{g0/0}. O roteador atua como gateway de cada segmento lógico por meio das subinterfaces \texttt{g0/0.10}, \texttt{g0/0.20}, \texttt{g0/0.30} e \texttt{g0/0.40}.

O provimento de concessões de endereço IP permanece desacoplado dos ativos de roteamento. Para viabilizar esse desacoplamento, as subinterfaces de acesso das estações de pesquisa e da rede sem fio recebem a instrução \texttt{ip helper-address 192.168.20.30}, que converte as transmissões de descoberta DHCP por difusão em mensagens direcionadas por unidifusão (\textit{unicast}) ao servidor autônomo \texttt{SRV-DHCP}.

Os enlaces de trânsito estruturam-se em três níveis hierárquicos:
\begin{enumerate}
    \item \textbf{Segmento WAN ISP-Borda (\texttt{200.17.0.0/30})}: Interconecta a interface \texttt{g0/0} do roteador do provedor \texttt{ISP-UFS} (\texttt{200.17.0.1}) à interface WAN \texttt{g0/0/1} do roteador corporativo \texttt{router\_border} (\texttt{200.17.0.2});
    \item \textbf{Segmento WAN-Firewall (\texttt{10.0.0.0/30})}: Interconecta a interface interna \texttt{g0/0/0} do Cisco ISR 4321 (\texttt{10.0.0.1}) à interface externa \texttt{g1/1} (\textit{outside}) do Cisco ASA 5506-X (\texttt{10.0.0.2});
    \item \textbf{Segmento Firewall-Núcleo (\texttt{10.0.1.0/30})}: Interconecta a interface interna \texttt{g1/2} (\textit{inside}) do Cisco ASA 5506-X (\texttt{10.0.1.1}) à interface de núcleo \texttt{g0/1} do Cisco 2911 (\texttt{10.0.1.2}).
\end{enumerate}

O enlace inter-edifícios é estabelecido entre a porta \texttt{GigabitEthernet0/2} do switch central \texttt{SW-CT12} e a porta \texttt{GigabitEthernet0/2} do switch local \texttt{SW-CT09}, operando em modo tronco 802.1Q nativo para transportar o tráfego segregado de todas as VLANs sobre o canal óptico monomodo.

\section{Políticas de Segurança, Inspeção de Estado e Filtragem L3/L4}

A proteção da infraestrutura é executada de forma coordenada em duas barreiras de filtragem:
\begin{enumerate}
    \item \textbf{Segurança Perimetral e Inspeção de Estado (Cisco ASA 5506-X)}: Bloqueio estrito de fluxos iniciados na zona não confiável (\textit{outside} para \textit{inside}) sem correspondência prévia na tabela de estados de conexão do \textit{firewall}. O tráfego interno originado na LAN dispõe de saída transparente, com retorno garantido pela inspeção MPF dos protocolos essenciais (ICMP, HTTP e DNS);
    \item \textbf{Isolamento Inter-VLAN e Controle Granular (Cisco 2911)}: Listas de controle de acesso estendidas filtram as conexões entre as redes virtuais. As estações de pesquisa e dispositivos Wi-Fi acessam exclusivamente as portas 80 e 443 do servidor \texttt{SRV-NEXTGIS} (\texttt{192.168.20.10}), enquanto a porta nativa do banco de dados PostgreSQL (TCP 5432) e os endereços da VLAN de gerência (VLAN 10) permanecem bloqueados para nós de acesso.
\end{enumerate}

A Listagem~\ref{lst:cisco_acl} detalha a ACL estendida aplicada na subinterface da rede de pesquisa, contendo a autorização explícita de pacotes BOOTP/DHCP na camada de acesso para viabilizar a negociação dinâmica com o servidor \texttt{SRV-DHCP}.

\begin{lstlisting}[language=bash,caption={Regras de ACL estendida aplicadas no gateway da rede de pesquisa.},label={lst:cisco_acl}]
ip access-list extended ACL_PESQUISA_IN
 permit udp any any eq bootpc
 permit udp any any eq bootps
 permit tcp 192.168.30.0 0.0.0.255 host 192.168.20.10 eq 80
 permit tcp 192.168.30.0 0.0.0.255 host 192.168.20.10 eq 443
 deny   tcp 192.168.30.0 0.0.0.255 host 192.168.20.20 eq 5432 log
 deny   ip 192.168.30.0 0.0.0.255 192.168.10.0 0.0.0.255
 permit ip 192.168.30.0 0.0.0.255 any
interface GigabitEthernet0/0.30
 ip helper-address 192.168.20.30
 ip access-group ACL_PESQUISA_IN in
\end{lstlisting}
